% Part: incompleteness
% Chapter: incompleteness-provability
% Section: godels-paper

\documentclass[../../../include/open-logic-section]{subfiles}

\begin{document}

% SEGMENT OLP-0317-S01

\olfileid{inc}{inp}{gop}

\olsection{கோடலின் முதற் கட்டுரையுடன் ஒப்பீடு}

கோடலின் 1931 ஆம் ஆண்டுக் கட்டுரையுடன் சிறிது நேரம் செலவிடுவது
பயனுள்ளது. நாம் இப்போது விவாதித்த எண்ணங்களை அதன் முன்னுரை
சுருக்கமாகக் கூறுகிறது. கட்டுரையை மேலோட்டமாகப் படித்தாலும் ஒவ்வொரு
கட்டத்திலும் என்ன நடக்கிறது என்பது எளிதில் தெரியும்: முதலில் கோடல்
$P$ என்ற முறைசார் அமைப்பை (தொடரியல், அடிகோள்கள், நிறுவல் விதிகள்)
வரையறுக்கிறார்; பின்னர் முதன்மை மீள்வரையறைச் சார்புகளையும் உறவுகளையும்
வரையறுக்கிறார்; அதன் பிறகு $x B y$ முதன்மை மீள்வரையறை என்பதை
காட்டி, முதன்மை மீள்வரையறைச் சார்புகளும் உறவுகளும் $\Th{P}$ இல்
பிரதிநிதித்துவப்படுத்தப்படுகின்றன என்று வாதிடுகிறார். பின்னர் மேலே
கண்ட முழுமையின்மைத் தேற்றத்தை நிறுவுகிறார். மூன்றாம் பிரிவில்,
நிறுவிக்க முடியாத கூற்றை எண்கணிதத்தின் மொழியிலுள்ள ஒரு வாக்கியமாக
எடுக்கலாம் என்பதை அவர் காட்டுகிறார். இதுவே பீட்டா-துணைத்தேற்றத்தின்
தோற்றம்; $\Th{Q}$ இல் மீள்வரையறைச் சார்புகள் பிரதிநிதித்துவப்படுகின்றன
என்பதைக் காட்டும்போது தொடர்களைக் கையாள நாம் இதையே, $\beta$
துணைத்தேற்றத்தில், பயன்படுத்தினோம்.
போதுமான குறைந்தபட்ச அடிகோள்களின் கணத்தை கோடல் தனியாகப் பிரித்துக்
காட்டவில்லை; ஆனால் $\Th{Q}$ இந்தப் பணியைச் செய்யும் என்பதை இப்போது
அறிவோம். இறுதியாக, நான்காம் பிரிவில் இரண்டாம் முழுமையின்மைத்
தேற்றத்திற்கான ஒரு நிறுவல் வரைவை அவர் தருகிறார்.

\end{document}
